\section{The question being continued}
\label{sec:context}

The current ProveIt manuscript treats special polylogarithmic values,
cyclotomic relations, exact certificates and Herglotz arithmetic in one
framework. Its Herglotz chapter proves a complete formal classification
of $F(p/q)-F(1)$ and of $J(2/q)$. Its discovery chapter explicitly leaves
the general $J(p/q)$ symbol problem as a further direction
\cite{Repo}. That is the problem addressed here. The Gaussian reductions,
the $S_4$ certificate and the other completed continuations are not
reclaimed as new work.

For $x>0$, use the standard normalization
\begin{equation}\label{eq:def-FJ}
 F(x)=\sum_{n\ge1}\frac{\psi(nx)-\log(nx)}{n},
 \qquad J(x)=\int_0^1\frac{\log(1+t^x)}{1+t}\dd t.
\end{equation}
Radchenko--Zagier establish the finite cyclotomic dilogarithm formula for
rational $F$ and the identity
\begin{equation}\label{eq:J-F}
 J(x)=F(2x)-2F(x)+F(x/2)+\frac{\pi^2}{12x}.
\end{equation}
Their modified Herglotz function has a Hecke relation that will supply
explicit evaluations in Section~\ref{sec:identities} \cite{RZ}.
The broader two-parameter Herglotz--Zagier--Novikov function and its
specializations are studied by Choie--Kumar \cite{CK}.

There are two different questions. An analytic evaluation is an equality
of numbers with specified branches. A formal reduction asks whether the
finite dilogarithm expression admits reduction by the rationalized
five-term calculus. Section~\ref{sec:symbols} defines the latter exactly.
The classification below is complete for that calculus. It does not
classify every conceivable numerical equality among these periods.

\subsection{Results and their dependencies}

The new algebraic step is not another numerical rank computation. At an
odd conductor, doubling the root set gives an exact three-term
\emph{conductor} identity for exterior symbols. At an even conductor,
a symbol in the new dyadic layer cannot be absorbed by symbols from
the preceding field unless it is zero. These facts reduce the boundary
of $J(p/q)$ to two sparse finite-group calculations.

The resulting classification has a particularly simple shape. For
mixed parity, logarithmic and rational-dilogarithmic reducibility
coincide. For odd numerator and denominator, the only nontrivial
rational-dilogarithmic cases are the ratios between $1$, $3$ and $5$.
The condition modulo $2e$, not merely modulo $e$, is essential.

A Vieta descent then parametrizes all logarithmic ratios using the
recurrences $w_{j+1}=c w_j\pm w_{j-1}$ with $c$ even. This gives both
an exact enumerator and a sparse height law. Finally, a separate analytic
argument yields a uniform formula for neighboring ratios and evaluates
the exceptional value $J(3/5)$.
These analytic identities carry their full logarithmic and $\pi^2$
terms; they are not merely boundary certificates.

\subsection{Scope of the audit and claims of novelty}

The reference snapshot is \snapshot. The main source file is
\nolinkurl{chapters/09-herglotz.tex}, whose Git blob is
\texttt{b372a62797ae9161ddcfd9f4edd036ecfd87df36}. The dependency and
open-question checks also included
\nolinkurl{chapters/10-discovery.tex} and the existing
\nolinkurl{herglotz-cyclotomic-obstructions} report.

Statements described as new are continuations relative to this reviewed
snapshot, not claims of exhaustive bibliographic priority. Classical
functional equations, the existing all-conductor symbol theorem, and
rationalized algebraic $K$-theory facts are attributed as inputs. Ordinary
proofs are supplied for the new results. The accompanying finite tests
are implementation checks, not substitutes for these proofs.

Section~\ref{sec:audit} gives one precise correction to a statement in
the cited preprint. It does not purport to audit every assertion of that
paper or of the complete ProveIt manuscript. In particular, it is not a
claim that the manuscript's independently derived boundary formula is
incorrect.
